\textbf{Domain randomisation.} Tobin et al.\ randomise rendering to transfer a visual detector to a real robot \cite{tobin2017domain}; Peng et al.\ randomise dynamics to transfer a manipulation policy \cite{peng2018sim}. Xie et al.\ study sim-to-real transfer for a quadruped and find that transfer is possible without dynamics randomisation in their case study, which is a reason to ask what randomisation costs \cite{xie2021dynamics}. OpenAI et al.\ introduce Automatic Domain Randomization, which expands the ranges when performance at the range boundary is good enough \cite{openai2019solving}; our success-gated curriculum is a much simpler relative of that idea, and we did not reimplement ADR itself. Alghonaim and Johns benchmark visual randomisation choices \cite{alghonaim2020benchmarking}. Chen et al.\ give a theoretical account of when DR supports sim-to-real transfer \cite{chen2021understanding}.

\textbf{Robust and adaptive training.} EPOpt trains on a model ensemble and focuses on a worst-case fraction of trajectories \cite{rajeswaran2016epopt}. Yu et al.\ combine randomisation with online system identification \cite{yu2017preparing}. Chebotar et al.\ and DROPO adapt the randomisation distribution using real-world data \cite{chebotar2019closing,tiboni2022dropo}.

\textbf{Curricula.} Automatic curriculum learning for deep RL is surveyed in \cite{portelas2020automatic}. Our curriculum is the simplest member of that family: a threshold on training success widens a fixed-shape range. We reimplemented all baselines in a small CartPole setting; none of the cited methods was run, and our results say nothing about their performance.
